We close on the four contributions in the same order they appear in \S\ref{sec:intro}. \textbf{C1:} A single RoBERTa-large fine-tuning pipeline, with only per-cell variation in input fields and loss weights, populated all six (subtask, track) cells competitively. \textbf{C2:} Synthetic contrastive augmentation gave consistent small gains (+0.014 to +0.087 across cells) and pseudo-labelling unlocked the enhanced track for ST2 and ST3, but a pseudo recipe lost ground on ST1 enhanced and a rare-class synthetic top-up lost ground on ST3 base; we reported both kinds of result in the same tables. \textbf{C3:} A grid covering DeBERTa-v3, ModernBERT, and four Qwen variants up to 32B parameters showed no validated benefit over the encoder at this corpus scale; the 5-fold standard deviation on the strongest LLM comparator leaves that comparison underpowered rather than closed. \textbf{C4:} The enhanced-track headline numbers should be read as upper bounds because the rewritten fields carry signal that our string-mask sanity check could not rule out at the semantic level, and the practical-external (\textsf{pe}) ST3 class continued to limit macro-F1 after targeted augmentation, threshold calibration, and a basis-axis meta-learner. Touch\'{e} continues to be a productive platform for benchmarking argument-aware NLP, and we look forward to further contributions in future editions.
